\chapter{Glossary of Terms}

This glossary provides definitions for key terms used throughout this treatise. Terms are organized alphabetically for reference.

\section*{A}

\textbf{Agentic AI} --- AI systems capable of autonomous action, tool use, and multi-step reasoning without constant human guidance. Distinguished from chat-based AI by the ability to execute plans independently.

\textbf{Approval Point} --- A moment in an agentic workflow where human review and authorization is required before the AI proceeds. Key locus for DCF application in modern systems.

\textbf{Autonomy Spectrum} --- The range of AI independence from fully supervised (every action requires approval) to fully autonomous (AI acts independently). DCF principles apply differently at each point.

\textbf{Anticipatory Calibration}\index{Anticipatory Calibration} --- The practice of forming a hypothesis about what AI will produce before prompting, then comparing expectation to reality. Builds accurate mental models of AI capabilities over time. See Chapter 11.

\section*{B}

\textbf{Background Agents}\index{Background Agents} --- Agents spawned via the Task tool that run asynchronously while you continue other work. Trade engagement for efficiency---use for exploration, engage synchronously for judgment calls.

\section*{C}

\textbf{Chain-of-Thought (CoT)} --- A prompting technique where the AI is encouraged to show its reasoning process step-by-step, improving performance on complex tasks. Related to but distinct from DCF's emphasis on human-AI dialogue.

\textbf{Checkpoint Engagement} --- The DCF principle of applying Socratic questioning at strategic approval points rather than monitoring every AI action. Enables efficient oversight of autonomous systems.

\textbf{CLAUDE.md} --- A configuration file in Claude Code that provides persistent context, preferences, and instructions. Functions as externalized cognitive infrastructure in DCF.

\textbf{Co-creation} --- The DCF principle that collaborative human-AI work produces capabilities that neither possesses alone. Distinguishes partnership from mere tool use.

\textbf{Consequences (Tracing Implications)}\index{Consequences} --- Questioning that explores the logical implications and outcomes of ideas. In DCF: ``What are the consequences? What if you're wrong?'' Reveals whether an idea actually works in practice and surfaces unintended effects.

\textbf{Cognitive Augmentation} --- Enhancement of human thinking capacity through AI partnership. Distinguished from automation (AI replacing human thought) and from mere information access.

\textbf{Cognitive Infrastructure} --- Systems, documents, and configurations that shape and support thinking. In DCF, includes CLAUDE.md files, memory systems, and documentation architecture.

\textbf{Cognitive Scaffolding} --- Temporary AI support that builds human capability over time. Effective scaffolding leads to eventual independence, not permanent dependence.

\textbf{Compounding (Cognitive)}\index{Compounding} --- The phenomenon where AI-assisted learning produces accelerating returns: better questions lead to better responses, which deepen understanding, which enables even better questions. Occurs through the reinforcement loop of articulation, challenge, synthesis, and internalization. Distinguishes learning-stance practitioners from extraction-mode users.

\textbf{Context Budgeting}\index{Context Budgeting} --- The practice of managing what information exists in an AI conversation's context. Every context has explicit limits (token windows) and implicit limits (attention degradation). Effective budgeting distinguishes what must be pre-loaded from what can be retrieved on-demand, and identifies what can be safely dropped without losing capability.

\textbf{Context Engineering}\index{Context Engineering} --- The strategic layer of DCF concerned with information architecture: what context to construct before dialogue begins. Complements the Socratic Toolkit (tactical: how to reason) and mode selection (operational: what questions to ask). Core principles include relevance over completeness, fresh over stale, structured over sprawling, and explicit over implicit.

\textbf{Context Window} --- The maximum amount of text an LLM can process in a single interaction. Determines the practical limits of prompt chaining and conversation depth.

\textbf{Critical Rationalism}\index{Critical Rationalism} --- Karl Popper's philosophical stance that all knowledge is provisional, conjectural, and open to revision through criticism. In DCF, provides the foundation for treating AI outputs as hypotheses to be tested rather than truths to be accepted.

\section*{D}

\textbf{DCF (Dialectical Cognition Framework)} --- The integrated methodology for human-AI collaboration described in this treatise. Combines Socratic dialogue, recursive refinement, and metacognitive awareness.

\textbf{Dialectic} --- The process of arriving at truth through exchange of opposing ideas. In DCF, refers to the productive tension between human and AI perspectives.

\textbf{Double-Loop Learning}\index{Double-Loop Learning} --- Chris Argyris's concept distinguishing learning that adjusts actions within existing assumptions (single-loop) from learning that questions the assumptions themselves (double-loop). DCF's metacognitive emphasis enables double-loop learning by surfacing and examining underlying mental models.

\section*{E}

\textbf{Epistemic Humility} --- Recognition of the limits of one's knowledge. In DCF, applies to both human and AI: neither should claim certainty beyond their evidence.

\textbf{Evidence (Probing Reasons)}\index{Evidence} --- Questioning that examines the factual foundation of statements. In DCF: ``How do you know? What evidence supports this?'' Many beliefs rest on untested assumptions or weakly understood foundations---this operation surfaces them.

\textbf{Escape Paths}\index{Escape Paths} --- Questions or conditions embedded in prompts that help AI determine when it has produced sufficient output. Rather than relying solely on human judgment to terminate iteration, escape paths build exit conditions directly into prompt design. See ``Escape Paths: Designing Prompts with Built-in Exits.''

\textbf{Extended Mind Thesis} --- The philosophical position that cognitive processes can extend beyond the brain into the environment. Provides theoretical grounding for treating AI as cognitive augmentation.

\section*{F}

\textbf{Falsificationism}\index{Falsificationism} --- Karl Popper's principle that scientific hypotheses gain credibility not from confirmation but from surviving serious attempts at refutation. In DCF, grounds the practice of actively challenging AI outputs rather than seeking confirmation.

\textbf{Fusion of Horizons}\index{Fusion of Horizons} --- Hans-Georg Gadamer's concept (\emph{Horizontverschmelzung}) that understanding emerges through dialogue when different perspectives merge to create shared meaning. In DCF, explains why human-AI collaboration is partnership: genuine understanding requires dialogue that changes both parties.

\section*{H}

\textbf{Hallucination} --- AI generation of plausible but factually incorrect information. A key limitation requiring human verification and Socratic questioning.

\textbf{Hooks}\index{Hooks} --- Shell commands that execute automatically before or after Claude Code tool calls. Enable automated DCF checkpoints---for example, running tests after every file edit.

\textbf{Human-in-the-Loop (HITL)} --- System design pattern where human judgment is required at critical decision points. DCF applies to optimizing these interaction points.

\section*{L}

\textbf{Language as Infrastructure} --- The DCF principle that well-crafted documentation and prompts function as cognitive architecture, shaping what kinds of thinking become possible.

\textbf{LLM (Large Language Model)} --- Neural networks trained on massive text corpora that generate human-like text. The AI systems to which DCF primarily applies.

\section*{M}

\textbf{MCP (Model Context Protocol)} --- A standard for connecting LLMs to external tools and data sources. Extends AI capabilities beyond text generation.

\textbf{Memory System} --- Persistent storage mechanisms (like memories in Claude Code) that allow context to persist across sessions. Part of cognitive infrastructure in DCF.

\textbf{Meta-Metacognition} --- Thinking about thinking about thinking. In DCF, the awareness of one's own metacognitive processes, such as noticing that you're noticing your reasoning.

\textbf{Metacognition} --- Awareness and analysis of one's own thinking processes. A core DCF principle: reflect on process, not just outcome.

\textbf{Meta-Question}\index{Meta-Question} --- The practice of asking ``What question should I be asking?'' rather than seeking answers directly. Often the most powerful Socratic move, as users are frequently stuck not from lacking answers but from asking the wrong question. Elevated to a core principle in DCF.

\section*{O}

\textbf{Operationalization} --- Converting abstract insights into concrete, actionable artifacts (documents, configurations, workflows). The move from understanding to implementation.

\section*{P}

\textbf{Pattern Capture}\index{Pattern Capture} --- The practice of recognizing effective prompting patterns and codifying them as reusable Claude Code skills. When you discover an approach that works well for a category of tasks, extract its transferable essence and document it. See Section~\ref{sec:pattern-to-skill}.

\textbf{Plan Mode} --- A feature in agentic systems where the AI develops a detailed plan before execution, allowing human review of the approach. Natural DCF application point.

\textbf{Prompt} --- Input text given to an LLM to elicit a response. Quality of prompts significantly affects quality of outputs.

\textbf{Prompt Chaining} --- Connecting multiple prompts in sequence, where output from one becomes input to the next. Enables complex reasoning that exceeds single-prompt capability.

\section*{R}

\textbf{Recursive Refinement} --- The DCF principle of iteratively improving outputs until they converge on clarity. Each iteration's output becomes the next iteration's input.

\textbf{Reinforcement Loop (Cognitive)} --- The four-stage cycle through which AI dialogue builds durable capability: (1) Articulation---externalizing thinking; (2) Challenge---AI responses that force active processing; (3) Synthesis---integrating new with existing understanding; (4) Internalization---patterns becoming automatic. The mechanism behind cognitive compounding.

\textbf{Reinvention Addiction}\index{Reinvention Addiction} --- The anti-pattern of repeatedly solving the same type of problem from scratch instead of capturing effective patterns as reusable skills. Treats each session as isolated rather than building cumulative capability. See Anti-Pattern 12 in Chapter~\ref{ch:failure-modes}.

\textbf{Research Dialogue} --- Using AI for exploration and understanding rather than task completion. Distinct from extraction: produces understanding rather than artifacts, explores multiple paths, and expands capability rather than leaving it unchanged. Investment that compounds over time.

\textbf{ReAct (Reasoning + Acting)} --- An AI paradigm combining reasoning traces with action execution. Enables LLMs to interact with external tools and environments.

\section*{S}

\textbf{Session Lifecycle}\index{Session Lifecycle} --- The practice of managing conversation context over time. Includes knowing when to start fresh sessions vs.\ continue existing ones, and using \texttt{/dcf compact} to capture state before context limits are reached.

\textbf{Socratic Dialogue} --- A method of inquiry through questioning rather than assertion. In DCF, the primary mode of productive human-AI interaction.

\textbf{Socratic Fatigue}\index{Socratic Fatigue} --- Exhaustion from over-application of Socratic questioning, occurring when the collaborative energy of productive dialogue masks diminishing returns. The human tendency toward over-perfection when engaged in recursive refinement. Counteracted by designing prompts with escape paths and exit hooks.

\textbf{Socratic Prompting} --- The application of Socratic questioning techniques to LLM interaction. Asking ``why'' and ``how do you know'' rather than simply accepting outputs.

\textbf{Sycophancy} --- The tendency of LLMs to agree with users rather than provide accurate information. A limitation that Socratic questioning helps counteract.

\section*{T}

\textbf{Thinking Mirror} --- The core DCF metaphor: LLMs reflect and transform human thought. Quality of reflection depends on quality of input; distortions are possible but manageable.

\textbf{Token} --- The basic unit of text that LLMs process. Words are typically split into multiple tokens. Token limits constrain context windows.

\textbf{Tool Use} --- The ability of LLMs to invoke external functions (file access, web search, code execution). Expands AI capability beyond pure text generation.

\section*{V}

\textbf{Verification} --- The process of checking AI outputs for accuracy and appropriateness. Essential complement to Socratic dialogue in DCF.

\textbf{Via Negativa}\index{Via Negativa} --- A Socratic operation that defines by exclusion: ``What should we explicitly NOT do?'' Drawing from the apophatic tradition, this technique activates different reasoning than positive guidance alone. Explicitly stating what to avoid surfaces failure modes and anti-goals that would otherwise remain implicit. Particularly valuable for the \texttt{constrain} mode and before generative work where rework is expensive.

\section*{W}

\textbf{Workflow Composition}\index{Workflow Composition} --- The practice of chaining multiple DCF skill modes into sequences appropriate for common scenarios. Skills remain atomic (invoked individually), but recommended sequences provide guidance for multi-phase work. Examples: ``New Project'' (\texttt{onboard} $\rightarrow$ \texttt{architect} $\rightarrow$ \texttt{premortem}), ``Decision Point'' (\texttt{tradeoffs} $\rightarrow$ \texttt{challenge} $\rightarrow$ \texttt{decide}), ``Complex Task'' (\texttt{constrain} $\rightarrow$ \texttt{decompose} $\rightarrow$ \texttt{architect}), ``High-Stakes Decision'' (\texttt{assumptions} $\rightarrow$ \texttt{verify} $\rightarrow$ \texttt{challenge} $\rightarrow$ \texttt{decide}). Each transition is a checkpoint where users engage fully before proceeding.

\section*{Z}

\textbf{Zone of Proximal Development (ZPD)} --- Vygotsky's concept of tasks achievable with assistance but not independently. AI can expand the ZPD for human learners.

% ============================================================================
% EXERCISES
% ============================================================================

